Let us briefly recall the definition of the abovementioned isomorphisms in degree $p=0$.
According to the analogue of \cite[Lemma B.7]{GueWillYu_DyncomplKtheory} in even degree\footnote{The cited lemma only considers the odd case, but the even case can be proven completely analogously.}, every class in $\KK_*^\Gamma(\Cz(X),D)$ can be represented by a cycle of the form $\left(\HilbertMod_D\oplus\HilbertMod_D,\left(\begin{smallmatrix}0&U^*\\U&0\end{smallmatrix}\right)\!,\epsilon\oplus\epsilon,\pi\oplus\pi\right)$ with $U\in \PseudoLoc_\Gamma(X;D)$.
Then $U$ is a unitary modulo $\Roe_\Gamma(X;D)$ and the Paschke duality isomorphism is defined by mapping this cycle to the class represented by $U$. The second isomorphism then takes it to the class represented by a function $t\mapsto U_t$ which represents a unitary in $\PseudoLocLoc(X;D)/\Loc(X;D)$ and such that $U_0$ agrees with $U$ modulo $\Roe_\Gamma(X;D)$. Such a function can be constructed from the constant function $U$ by applying the procedure from the proof of \cite[Lemma~B.15]{GueWillYu_DyncomplKtheory}.
The third isomorphism then maps it to the class represented by the difference of the projection valued functions $t\mapsto P_t$ and $t\mapsto Q$ defined by
\[
P_t\coloneqq \begin{pmatrix}U_tU_t^*&U_t(\id-U_t^*U_t)^{\frac12} \\U_t^*(\id-U_tU_t^*)^{\frac12}&\id-U_t^*U_t\end{pmatrix}\qquad\text{and}\qquad Q=\begin{pmatrix}\id&0\\0&0\end{pmatrix}\!,
\]
cf.\ \cite[Proposition~4.8.10]{HigRoe}.

In order to prove the claim in even degrees, it now suffices to show that $\Delta$ composed with these isomorphisms is the canonical natural isomorphism $\KK_0^\Gamma(\Cz(X),D)\xrightarrow{\cong}\EE_0^\Gamma(\Cz(X),D)$.
The latter is defined via \cite[Theorem~2.2]{Thomsen_univKK} and hence we have to verify two properties: That the composition is natural in $D$ and that it maps $1_{\Cz(X)}\in\KK_0^\Gamma(\Cz(X),\Cz(X))$ to $1_{\Cz(X)}\in\EE_0^\Gamma(\Cz(X),\Cz(X))$. Naturality is clear from the description of the three isomorphisms above and the definition of~$\Delta$.

For the verification of the other property, recall that $1_{\Cz(X)}\in\KK_0^\Gamma(\Cz(X),\Cz(X))$ is represented by the cycle $(\Cz(X)\oplus 0,0,\varepsilon'\oplus 0,\pi'\oplus 0)$, where $\varepsilon'$ and $\pi'$ denote the canonical representations of $\Gamma$ and $\Cz(X)$. After adding the degenerate cycle $\left(\HilbertMod_{\Cz(X)}\oplus\HilbertMod_{\Cz(X)}, \left(\begin{smallmatrix}0&\id\\\id&0\end{smallmatrix}\right)\!,\epsilon\oplus\epsilon,\pi\oplus\pi\right)$, we see that it is also represented by the cycle $\left(\HilbertMod_{\Cz(X)}\oplus\HilbertMod_{\Cz(X)}, \left(\begin{smallmatrix}0&U^*\\U&0\end{smallmatrix}\right)\!,\epsilon\oplus\epsilon,\pi\oplus\pi\right)$ for a~projection $U\colon \HilbertMod_{\Cz(X)}\cong\Cz(X)\oplus\HilbertMod_{\Cz(X)}\to\HilbertMod_{\Cz(X)}$ onto the second summand. Since $U$ is a module homomorphism, it has propagation zero and $U\in \PseudoLocLoc(X;\Cz(X))$. Furthermore we have $UU^*=\id$ and $p\coloneqq \id-U^*U$ is a projection onto an isomorphically embedded copy of $\Cz(X)$ in $\HilbertMod_{\Cz(X)}$.
For the calculation of the image of this class in $\EE_0^\Gamma(\Cz(X),D)$ we can choose the constant function $t\mapsto U_t\coloneqq U$. We end up with the element that is represented by the $*$-homomorphism $\Cz(X)\to \Kom(\HilbertMod_{\Cz(X)})$, $f\mapsto \pi(f)p$, and this is exactly $1_{\Cz(X)}$.

The isomorphism in odd degrees now follows from the one in even degrees by replacing $D$ with the suspension $D\otimes \Cz(\R)$.
\end{proof}
